\section{Results}
\label{sec:results}

\subsection{Band-order sensitivity and matched controls}

\begin{table}[htbp]
\centering
\caption{Three-seed mean performance under both plausible HR-GLDD orders.
BF1 is boundary F1 with one-pixel Chebyshev tolerance. Neither order is
declared authoritative.}
\label{tab:controlled}
\footnotesize
\setlength{\tabcolsep}{3.5pt}
\begin{tabular}{@{}lrrrr@{}}
\toprule
Configuration & RGBN F1 & BGRN F1 & RGBN BF1 & BGRN BF1 \\
\midrule
U-Net/base & 66.92 & 66.92 & 49.62 & 49.62 \\
ResU-Net/base & 68.55 & 68.55 & 52.93 & 52.93 \\
Attn-zero/base & 68.88 & 68.88 & 53.17 & 53.17 \\
Attn-raw/base & 69.05 & 69.13 & 52.92 & 53.15 \\
Attn-index/base & 69.04 & 68.92 & 52.76 & 52.83 \\
Attn-zero/plain-boundary & 69.99 & 69.99 & 56.04 & 56.04 \\
Attn-raw/plain-boundary & 69.88 & 70.19 & 55.85 & 56.15 \\
Attn-index/plain-boundary & 69.98 & 70.12 & 55.96 & 56.20 \\
Attn-index/index-boundary & 69.58 & 69.78 & 56.34 & 56.15 \\
\bottomrule
\end{tabular}

\end{table}

Table~\ref{tab:controlled} reports every configuration under both candidate
orders; Supplementary Table S3 gives detailed RGBN-candidate operating metrics.
With base loss, candidate-index attention reaches $69.04\%$ F1 under
RGBN and $68.92\%$ under BGRN. Raw candidate-Red/NIR edges reach $69.05\%$ and
$69.13\%$, while the capacity-matched zero control reaches $68.88\%$ under
both. Candidate-index input specificity is therefore unsupported regardless
of order.
All contrasts and envelopes below use unrounded run-level values; subtracting
the table's two-decimal cells can differ by $0.01$ point.

Plain boundary weighting raises F1 in every seed and control mode under both
orders. Mean F1 spans $69.88$--$70.19\%$ and boundary F1
$55.85$--$56.20\%$. Equal-mass index modulation changes F1 by $-0.41$ points
under RGBN and $-0.33$ under BGRN. Its boundary-F1 contrast changes from
$+0.38$ to $-0.05$ points, so spectral modulation is not order-robust. The
order-stable result is generic boundary emphasis, not spectral identity.

\subsection{Semantic-identifiability certificate}

The index-input-versus-raw F1 envelope is $[-0.21,-0.01]$ points and the
index-versus-plain-loss F1 envelope $[-0.41,-0.33]$, identifying negative
directions over the two admissible schemas. The corresponding boundary-F1
envelope $[-0.05,+0.38]$ crosses zero and is direction-unidentified. Generic
plain-boundary F1 has envelope $[+0.96,+1.13]$; all 18 individual effects lie
between $+0.42$ and $+1.62$ points. Supplementary Table~S4 is the generated
certificate.

\subsection{External analyses}

The post-hoc CAS sensitivity analysis fails three of four quantitative
criteria. Across three designated regions and seeds, plain boundary weighting
changes mean region-macro F1 by $-2.19$ points and boundary F1 by $-1.86$
points; seed F1 effects are $-10.42$, $+1.22$, and $+2.62$ points.
Supplementary Table~S1 reports every contrast. Because the protocol was first
committed after execution, this result receives no confirmation credit.

The internally precommitted LRD run fails all four frozen criteria, but is not
confirmation-grade. Maximum validation F1 is only $0.008$--$0.151$; protected
event-macro F1 spans $0.007$--$0.283$ for base and $0.009$--$0.132$ for
boundary weighting. Under selected thresholds, mean F1 changes by $-4.73$
points. The frozen per-crop boundary convention yields $-24.10$ points, but
crop-macro empty/empty scoring of zero yields $-0.15$, pooled selected-threshold
scoring $-4.92$, and pooled common-$0.5$ scoring $-0.55$. Development PH0003 and
protected PH0001/PH0004 also share a trigger
family, leaving four non-overlapping protected EIDs. This external evidence is
failed/indeterminate, not a robust transfer rejection. Supplementary Table~S2
and Method~S2 report the frozen and sensitivity results.

\subsection{Operating point and spatial degradation}

At threshold $0.5$, boundary-weighted residual-attention FPR spans
$2.47$--$2.50\%$ across orders. Validation-selected nearest recall breakpoints
match exactly in 53 of 54 order-specific runs; the RGBN
index/plain-boundary seed-42 run differs by $2.11\times10^{-6}$. Combined
matched FPR and precision ranges across the seven residual-attention
configurations are $1.97$--$2.01\%$ and
$78.41$--$78.74\%$. These differences do not support index-specific
false-positive suppression.

Controlled $128\rightarrow38\rightarrow128$ degradation lowers F1 by
$4.33$--$5.20$ points across both orders. Raw-edge plain-boundary F1 changes
from $69.88$ to $65.37\%$ under RGBN and from $70.19$ to $65.56\%$ under
BGRN. This identifies same-sensor spatial-information sensitivity, not
cross-sensor causation.
